% Supported v0.5 teaching-philosophy-statement example.
\documentclass[type=teaching-philosophy]{careerdossier-statement}
% examples/profiles/profile-academic.tex
%
% Shared profile metadata used by the academic-CV example and the corresponding
% industry-resume smoke fixture. Keeping it independent of document layout
% proves that both document families consume the same \CDossierSetup interface.

\CDossierSetup{
  name     = {Ada Lovelace},
  headline = {Researcher in Analytical Computing},
  email    = {ada.lovelace@example.com},
  location = {Toronto, Ontario, Canada},
  website  = {example.com/ada-lovelace},
  %% The bare Google Scholar identifier, expanded to
  %% scholar.google.com/citations?user=ada-example for both the displayed text
  %% and the link. Pasting the whole address works too and is left as written.
  %% See docs/API.md, "Accepted forms for the web-profile keys".
  scholar  = {ada-example},
  orcid    = {0000-0002-1825-0097},
  affiliation = {Example University},
}

\CDossierStatementSetup{
  subtitle={Learning as inquiry, practice, and reflection},
  application-context={Faculty application, Example University}
}
\begin{document}
\MakeCDossierStatementHeader

\section*{Learning as an Active Construction}

I understand learning as the process of building, testing, and revising a way
of thinking. Students develop durable knowledge when they use ideas to make
predictions, encounter the limits of an initial explanation, and refine it in
conversation with evidence and other people. My role is to design conditions in
which that work is possible: meaningful questions, clear expectations,
well-timed guidance, and enough intellectual safety for students to expose
uncertainty without mistaking comfort for the absence of rigor.

This philosophy makes transparency central to teaching. I explain why an
activity matters, what successful reasoning looks like, and how a course's
parts connect. Transparency does not remove challenge. It directs effort toward
the disciplinary problem rather than toward guessing what an instructor wants.
As students gain experience, I gradually transfer responsibility for planning,
evaluation, and revision to them.

\section*{From Belief to Practice}

I begin new topics with a problem that reveals the need for the idea we are
about to study. Students first articulate what they already know and where that
knowledge becomes insufficient. A concise explanation or demonstration then
gives them tools to revisit the problem. Practice alternates between individual
reasoning, paired comparison, and whole-group synthesis so that students must
both form and communicate a position.

In computational settings, I ask students to treat an output as the beginning
of interpretation rather than the end of a task. They annotate assumptions,
design a test that could expose an error, and explain how numerical or data
choices affect the conclusion. This routine develops technical judgment while
also showing that responsible work is collaborative: another person should be
able to inspect the reasoning and continue from it.

Examples and counterexamples are especially important in my teaching. A correct
example can demonstrate a procedure, but a carefully chosen failure reveals the
conditions under which an idea is meaningful. I invite students to propose
their own edge cases and compare them. The resulting discussion positions
students as contributors to the course's knowledge rather than recipients of a
finished account.

\section*{Assessment and Feedback as Learning}

I use assessment to make thinking observable and to guide the next stage of
learning. Frequent low-stakes tasks identify misconceptions while they are
still revisable. Larger assignments require students to connect conceptual
claims, methods, validation, and communication. Rubrics describe these
dimensions separately, and exemplars are discussed before submission so that
criteria become tools students can use rather than rules revealed afterward.

Feedback is most valuable when students must interpret and apply it. I reserve
time for revision, ask students to identify the change that most improved their
work, and use short conferences when written comments alone are unlikely to
resolve the issue. I also seek feedback on the course through anonymous prompts
and patterns in student work. Revising a lesson in response models the same
reflective practice I ask of students.

\section*{Inclusive and Continuing Practice}

An inclusive learning environment offers genuine access to demanding work. I
provide materials in usable formats, define unfamiliar conventions, vary modes
of participation, and structure collaboration so that speaking quickly is not
the only route to intellectual visibility. I learn students' names, invite them
to connect course questions with their prior knowledge, and examine whose
examples and applications are treated as central.

My teaching philosophy is therefore not a fixed declaration. It is a commitment
to inquiry into my own practice. At Example University, I would continue that
work with colleagues and students, contribute to evidence-informed curriculum
design, and create courses in which learners leave with both stronger technical
abilities and a clearer sense of how to direct their own learning.

\end{document}
